\chapter{Experimental Evaluation}
\label{chap:evaluation}

This chapter reports the empirical results collected from the supplied paper and poster resources. The evaluation compares standard \DINO-family encoders with \DeACL-tuned \DINO ViT-B/16 across classification, semantic segmentation, and depth estimation.

\section{Experimental Setup}

The standard encoders are \DINO ViT-B/16, \DINOv ViT-S/14, and \DINOv ViT-B/14. The robust encoder is \DINO ViT-B/16 after \DeACL fine-tuning. Downstream heads are trained on frozen representations so that the evaluation isolates encoder quality and robustness.

Classification follows linear probing on CIFAR10, CIFAR100, and STL10 \parencite{krizhevsky2009learning,coates2011analysis}. Semantic segmentation is evaluated on ADE20K, Cityscapes, and PASCAL VOC 2012 \parencite{ADE1,ADE2,Cordts2016cityScapes,Everingham10pascalVoc}. Depth estimation is evaluated on NYU-Depth v2 \parencite{Silberman2012NYU}. Hyperparameters are listed in \Cref{sec:hyperparameters}.

\section{Linear Classification}

\Cref{tab:classification} shows that standard encoders have high clean accuracy but almost no adversarial accuracy. For example, \DINOv ViT-B/14 reaches 0.98 clean accuracy on CIFAR10 and 0.99 on STL10, but task-specific \PGD reduces all standard rows to 0.00.

\begin{table}[t]
    \centering
    \small
    \resizebox{\textwidth}{!}{%
    \begin{tabular}{lllccc}
    \toprule
    Dataset & SSL framework & Encoder type & Clean accuracy $\uparrow$ & \EmbedAttack accuracy $\uparrow$ & \PGD accuracy $\uparrow$ \\
    \midrule
    CIFAR10 & \DINOv ViT-S/14 & Standard & 0.94 & 0.01 & 0.00 \\
    CIFAR10 & \DINOv ViT-B/14 & Standard & 0.98 & 0.04 & 0.00 \\
    CIFAR10 & \DINO ViT-B/16 & Standard & 0.94 & 0.01 & 0.00 \\
    CIFAR10 & \DINO ViT-B/16 & \DeACL & 0.91 & 0.73 & 0.02 \\
    \midrule
    CIFAR100 & \DINOv ViT-S/14 & Standard & 0.82 & 0.00 & 0.00 \\
    CIFAR100 & \DINOv ViT-B/14 & Standard & 0.86 & 0.00 & 0.00 \\
    CIFAR100 & \DINO ViT-B/16 & Standard & 0.76 & 0.00 & 0.00 \\
    CIFAR100 & \DINO ViT-B/16 & \DeACL & 0.72 & 0.55 & 0.03 \\
    \midrule
    STL10 & \DINOv ViT-S/14 & Standard & 0.98 & 0.06 & 0.00 \\
    STL10 & \DINOv ViT-B/14 & Standard & 0.99 & 0.20 & 0.00 \\
    STL10 & \DINO ViT-B/16 & Standard & 0.98 & 0.00 & 0.00 \\
    STL10 & \DINO ViT-B/16 & \DeACL & 0.97 & 0.83 & 0.23 \\
    \bottomrule
    \end{tabular}}
    \caption{Linear classification robustness. \DeACL strongly improves \EmbedAttack robustness but leaves task-specific \PGD accuracy low.}
    \label{tab:classification}
\end{table}

\DeACL changes the representation-attack results substantially. On CIFAR10, \EmbedAttack accuracy rises from 0.01 to 0.73 for \DINO ViT-B/16. On STL10 it rises from 0.00 to 0.83. However, the downstream \PGD attack remains strong: 0.02 on CIFAR10, 0.03 on CIFAR100, and 0.23 on STL10. The defense therefore helps most when the attack matches its representation-space training objective.

\section{Semantic Segmentation}

\Cref{tab:segmentation} reports mIoU. Clean performance shows the value of stronger foundation encoders: \DINOv reaches 0.65--0.68 on Cityscapes and 0.83 on PASCAL VOC 2012. Yet adversarial performance is near zero for standard encoders under both \EmbedAttack and \SegPGD.

\begin{table}[t]
    \centering
    \small
    \resizebox{\textwidth}{!}{%
    \begin{tabular}{lllccc}
    \toprule
    Dataset & SSL framework & Encoder type & Clean mIoU $\uparrow$ & \EmbedAttack mIoU $\uparrow$ & \SegPGD mIoU $\uparrow$ \\
    \midrule
    ADE20K & \DINOv ViT-S/14 & Standard & 0.42 & 0.01 & 0.01 \\
    ADE20K & \DINOv ViT-B/14 & Standard & 0.45 & 0.00 & 0.01 \\
    ADE20K & \DINO ViT-B/16 & Standard & 0.27 & 0.01 & 0.01 \\
    ADE20K & \DINO ViT-B/16 & \DeACL & 0.24 & 0.14 & 0.01 \\
    \midrule
    Cityscapes & \DINOv ViT-S/14 & Standard & 0.65 & 0.02 & 0.01 \\
    Cityscapes & \DINOv ViT-B/14 & Standard & 0.68 & 0.03 & 0.00 \\
    Cityscapes & \DINO ViT-B/16 & Standard & 0.45 & 0.06 & 0.06 \\
    Cityscapes & \DINO ViT-B/16 & \DeACL & 0.36 & 0.31 & 0.03 \\
    \midrule
    PASCAL VOC 2012 & \DINOv ViT-S/14 & Standard & 0.83 & 0.00 & 0.01 \\
    PASCAL VOC 2012 & \DINOv ViT-B/14 & Standard & 0.83 & 0.00 & 0.01 \\
    PASCAL VOC 2012 & \DINO ViT-B/16 & Standard & 0.56 & 0.06 & 0.00 \\
    PASCAL VOC 2012 & \DINO ViT-B/16 & \DeACL & 0.51 & 0.30 & 0.02 \\
    \bottomrule
    \end{tabular}}
    \caption{Semantic segmentation robustness. \DeACL improves \EmbedAttack mIoU, but \SegPGD remains effective.}
    \label{tab:segmentation}
\end{table}

The \DeACL pattern repeats. \EmbedAttack mIoU improves from 0.06 to 0.31 on Cityscapes and from 0.06 to 0.30 on PASCAL VOC 2012. But \SegPGD mIoU remains 0.03 and 0.02 respectively. Clean mIoU also drops compared with standard \DINO ViT-B/16. Dense prediction therefore exposes the weakness of treating representation robustness as task robustness.

\section{Depth Estimation}

Depth estimation uses RMSE, so lower values are better. \Cref{tab:depth} shows that \DINOv ViT-B/14 has the best clean RMSE, 0.46, while \DeACL has worse clean RMSE than standard \DINO ViT-B/16.

\begin{table}[t]
    \centering
    \small
    \resizebox{\textwidth}{!}{%
    \begin{tabular}{llccc}
    \toprule
    SSL framework & Encoder type & Clean RMSE $\downarrow$ & \EmbedAttack RMSE $\downarrow$ & \DepthPGD RMSE $\downarrow$ \\
    \midrule
    \DINOv ViT-S/14 & Standard & 0.49 & 1.54 & 2.60 \\
    \DINOv ViT-B/14 & Standard & 0.46 & 1.29 & 2.74 \\
    \DINO ViT-B/16 & Standard & 0.61 & 1.28 & 1.79 \\
    \DINO ViT-B/16 & \DeACL & 0.68 & 0.92 & 1.71 \\
    \bottomrule
    \end{tabular}}
    \caption{NYU-Depth v2 robustness. \DeACL reduces \EmbedAttack damage but barely changes \DepthPGD damage.}
    \label{tab:depth}
\end{table}

\EmbedAttack substantially degrades depth prediction even though it does not optimize RMSE directly. For standard \DINO ViT-B/16, RMSE rises from 0.61 to 1.28. \DeACL improves this attacked RMSE to 0.92, but \DepthPGD remains severe at 1.71. The task-specific attack is stronger because it directly optimizes the geometric output loss.

\section{Fine-Tuning Dynamics}

\Cref{fig:deacl-dynamics} visualizes the \DeACL fine-tuning curves from the supplied figures. Most gains against \EmbedAttack appear early, while task-specific robustness remains limited. This supports the interpretation that the fine-tuning objective is well matched to representation attacks but insufficient for downstream losses.

\begin{figure}[t]
    \centering
    \begin{minipage}[t]{0.32\linewidth}
        \centering
        \includegraphics[width=\linewidth]{classification_deacl.pdf}
        \caption*{Classification}
    \end{minipage}\hfill
    \begin{minipage}[t]{0.32\linewidth}
        \centering
        \includegraphics[width=\linewidth]{segmentation_deacl.pdf}
        \caption*{Segmentation}
    \end{minipage}\hfill
    \begin{minipage}[t]{0.32\linewidth}
        \centering
        \includegraphics[width=\linewidth]{depth_estimation_deacl.pdf}
        \caption*{Depth}
    \end{minipage}
    \caption{\DeACL fine-tuning dynamics across downstream tasks. The curves show that representation-attack robustness improves more clearly than task-specific robustness.}
    \label{fig:deacl-dynamics}
\end{figure}

\section{Summary of Findings}

The empirical findings are consistent across tasks.

\begin{itemize}
    \item Standard SSL encoders are clean-useful but adversarially brittle.
    \item \EmbedAttack can strongly damage downstream performance even without task labels.
    \item \DeACL improves robustness against \EmbedAttack, the attack surface it trains against.
    \item Task-specific attacks remain effective, especially for segmentation and depth.
    \item Robustness of a foundation encoder should be reported per task and per attack interface.
\end{itemize}
